% 张三慧第四版第一卷第10章《热力学第一定律》
% 原教材PDF物理页371--375；思考题10.1--10.9不计入章末习题。
% 所有图10.22--10.32的箭头为核验依据，候选稿待独立审核。
\section*{第10章\quad 热力学第一定律}
\subsection*{习题10.1—10.26}

\paragraph{习题 10.1}
一定质量理想气体经历图10.22给出的
\(p\)-\(V\) 路径，\(BC\) 为以 \(pV\) 为常数的
双曲线；状态点 \(A=(10\,\mathrm L,2\,\mathrm{atm})\)，
\(B=(20\,\mathrm L,2\,\mathrm{atm})\)，
\(C=(40\,\mathrm L,1\,\mathrm{atm})\)，
\(D=(20\,\mathrm L,1\,\mathrm{atm})\)。
（1）若 \(T_A=300\,\mathrm K\)，求 \(T_B,T_C,T_D\)；
（2）求从 \(A\) 到 \(D\) 按图示经过 \(B,C\) 的
总气体功；（3）在 \(V\)-\(T\) 图标出状态及路径方向。
\textit{原图A—B、B—C、C—D箭头须在高分辨率图片中核对。}

\paragraph{习题 10.2}
图10.23中一热力学系统从状态 \(a\)
沿 \(acb\) 到 \(b\)，吸热 \(560\,\mathrm J\)，
对外做功 \(356\,\mathrm J\)。
（1）若改沿 \(adb\) 路径，
对外做功 \(220\,\mathrm J\)，求吸热；
（2）若沿路径 \(ba\) 从 \(b\) 回 \(a\)，
外界对它做功 \(282\,\mathrm J\)，
求该过程放热还是吸热以及热量大小。

\paragraph{习题 10.3}
\(64\,\mathrm g\) 氧气从 \(0^\circ\mathrm C\)
升至 \(50^\circ\mathrm C\)，比较
（1）体积保持不变与（2）压强保持不变
两过程的吸热、内能变化和对外功。
氧气在此温区取双原子理想气体模型。

\paragraph{习题 10.4}
\(10\,\mathrm g\) 氮气在压强不变下
吸收 \(10^3\,\mathrm J\) 的热量。
初温 \(300\,\mathrm K\)，求最终温度。

\paragraph{习题 10.5}
一定量氢气在
\(4.00\times10^5\,\mathrm{Pa}\)
等压下由 \(0.0^\circ\mathrm C\)
升温至 \(50.0^\circ\mathrm C\)，
吸热 \(6.0\times10^4\,\mathrm J\)。
求（1）气体的摩尔数；（2）内能变化；
（3）对外功；（4）若定容升高同样温度，
吸热量多少。氢气取双原子理想气体。

\paragraph{习题 10.6}
对 \(n\) 摩尔理想气体，
分别从给定初态到终态经历
\emph{等体、等温、绝热、等压}四种
准静态过程，用初末状态参量、自由度和热容
表示气体对外做功 \(A\)、吸热 \(Q\) 与
内能变化 \(\Delta E\)，并填入原书比较表。
\textit{本稿约定 \(A>0\) 为气体对外功、
\(Q>0\) 为吸热。}

\paragraph{习题 10.7}
一摩尔氧气绝热地从
\(10\,\mathrm L\)、\(300\,\mathrm K\)
膨胀到 \(30\,\mathrm L\)。
求最终压强、温度、对外做功、
内能变化；取 \(\gamma=1.40\)。

\paragraph{习题 10.8}
\(3\,\mathrm{mol}\) 气体初态为
\(p=2\,\mathrm{atm}\)、\(V=40\,\mathrm L\)，
先绝热压缩到一半体积，
随后等温膨胀回原体积。
取 \(\gamma=1.40\)。
（1）求全过程最大压强和最高温度；
（2）求全过程吸热、对外功、内能变化；
（3）画出 \(p\)-\(V\) 路径。

\paragraph{习题 10.9}
图10.24一绝热气缸中：
最初气缸总体积 \(30\,\mathrm L\)，
其中一侧氮气质量 \(35\,\mathrm g\)，
占体积 \(20\,\mathrm L\)，
压强 \(2\,\mathrm{atm}\)；
其余 \(10\,\mathrm L\) 为真空。
打开隔板小孔使氮气自由扩散充满
整个初始 \(30\,\mathrm L\) 的气缸；
然后缓慢移动活塞使气体绝热膨胀至
总体积 \(50\,\mathrm L\)。
求最终压强、温度、
整个过程对外功与内能变化。
\textit{自由扩散和缓慢绝热膨胀是两个不同阶段；
使用主教材图10.24确认活塞与腔室初体积。}

\paragraph{习题 10.10}
标准状态下，氧气中声速
\(3.172\times10^2\,\mathrm{m/s}\)。
利用声速—气体热容比关系求
氧气的比热比 \(\gamma=C_p/C_V\)。

\paragraph{习题 10.11}
按\emph{大气绝热模型}，证明气压随高度 \(h\) 变化：
\[
 p=p_0\left(
 1-\frac{Mgh}{C_{p,m}T_0}
 \right)^{\gamma/(\gamma-1)}.
\]
\(p_0,T_0\) 为 \(h=0\) 的压强温度，
\(C_{p,m}\) 是气体摩尔定压热容，
\(M\) 是摩尔质量。

\paragraph{习题 10.12}
墨西哥海滩的人体炮弹表演：
炮筒直径 \(0.80\,\mathrm m\)，长度
\(4.0\,\mathrm m\)，炮筒底部燃烧气体
初容积 \(2.0\,\mathrm{m^3}\)，初压
\(2.7\,\mathrm{atm}\)；燃气绝热推动炮弹，
排出口环境为 \(1.0\,\mathrm{atm}\)，
比热比 \(\gamma=1.4\)。
气体推动时对外功的75%传给体重
\(70\,\mathrm{kg}\) 的表演者。
求出口速率。
\textit{忽略摩擦和大气阻力，排口体积需加炮筒截面积×长度。}

\paragraph{习题 10.13}
证明一定量气体绝热地经过节流装置后，
如果前态 \((p_1,V_1,E_1)\)，
后态 \((p_2,V_2,E_2)\)，则
\[
 E_1+p_1V_1=E_2+p_2V_2.
\]
说明系统焓 \(H=E+pV\)
在此理想节流过程中守恒。

\paragraph{习题 10.14}
一固定质量气体初态 \((T_0,p_0,V_0)\)。
在\emph{第一实验}中，电流定时定容加热到
\((T_1,p_1,V_0)\)；
在\emph{第二实验}中，取同一起始态，
电流加热时间、功率不变，但等压
升温到 \((T_2,p_0,V_1)\)。
证明
\[
 \gamma=\frac{(p_1-p_0)V_0}
 {p_0(V_1-V_0)}.
\]
\textit{原题“第二次加热”应视为同初始态的比较实验，
不能误连为从第一实验的终态继续等压加热。}

\paragraph{习题 10.15}
理想气体的多方过程定义为 \(pV^n=\mathrm{const}\)。
（1）说明 \(n=0,1,\gamma,\infty\) 所代表的
等压、等温、绝热、等容过程；
（2）推导气体对外功的表达式；
（3）求多方过程摩尔热容 \(C_m\)，
说明特殊指数时的热容。
\textit{对功使用本项目统一的“气体对外为正”符号。}

\paragraph{习题 10.16}
绝热容器总体积
\(40\,\mathrm L\)，隔板可自由无摩擦滑动，
把气体分成 \(A,B\) 两部分，
每部分初装 \(1\,\mathrm{mol}\) 氮气，
两侧初始压强
\(1.013\times10^5\,\mathrm{Pa}\)，
各占 \(20\,\mathrm L\)。
只加热B室，隔板移动，
A室气体绝热压缩至原体积的一半。
求（1）B室 \(V\)-\(T\) 路径方程；
（2）两室末温；（3）B室吸热量。
以 \(\gamma=1.4\)、容器刚性绝热处理。

\paragraph{习题 10.17}
以图10.28所示理想布雷顿循环（Brayton cycle）：
\(1\to2\) 绝热压缩，
\(2\to3\) 等压吸热，
\(3\to4\) 绝热膨胀做功，
\(4\to1\) 等压放热。
（1）证明
\(\eta=1-(T_4-T_1)/(T_3-T_2)\)；
（2）若压缩比
\(r_p=p_{\max}/p_{\min}\)，证明
\(\eta=1-r_p^{-(\gamma-1)/\gamma}\)，
\(\gamma=1.4,r_p=10\) 时求数值。

\paragraph{习题 10.18}
图10.29为理想柴油机循环：
\(a\to b\)、\(c\to d\) 绝热，
\(b\to c\) 等压吸热，
\(d\to a\) 等容放热。
以最大容积 \(V_1\)、压缩后体积 \(V_2\)
及等压燃烧后容积 \(V'_1\)，推导效率
\[
 \eta=1-
 \frac{(V'_1/V_2)^\gamma-1}
 {\gamma(V_1/V_2)^{\gamma-1}(V'_1/V_2-1)}.
\]

\paragraph{习题 10.19}
图10.30给出三热库可逆循环，
与温度为 \(T,T_2,T_1\) 的热库接触的
等温过程分别为 \(ab,cd,ef\)，
其余 \(bc,de,fa\) 为绝热。
在 \(cd\) 过程中吸收的热量
恰等于在 \(ef\) 过程中放出的热量。
系统为理想气体，求效率 \(\eta\)
以 \(T,T_2,T_1\) 表示。

\paragraph{习题 10.20}
两台卡诺热机串联：
第一台低温热库作为第二台的高温热库。
证明总效率
\[
 \eta=\eta_1+(1-\eta_1)\eta_2,
\]
并证明总效率与单一台直接在
最高、最低两热库间运行的卡诺热机相同。

\paragraph{习题 10.21}
利用表层海水 \(25^\circ\mathrm C\)
和深层海水 \(5^\circ\mathrm C\)
作理想热机的两热库。
（1）求理论卡诺效率；
（2）若机械功率为 \(1\,\mathrm{MW}\)，
以何速率排出废热；
（3）获得机械功及排出的废热均来源于
表层海水从25℃冷却到5℃释放的热量，
问每秒抽取表层海水的体积流率。
\textit{海水比热以水 \(4.18\,\mathrm{kJ/(kg\,K)}\)
近似，密度以 \(1000\,\mathrm{kg/m^3}\) 估算。}

\paragraph{习题 10.22}
冰箱冷冻室内温度 \(-10^\circ\mathrm C\)，
室温 \(15^\circ\mathrm C\)。
按理想卡诺制冷循环，
若冰箱每消耗 \(10^3\,\mathrm J\) 电功，
可从冷冻室吸收多少热量？

\paragraph{习题 10.23}
室外 \(32^\circ\mathrm C\)，
用空调维持室内 \(21^\circ\mathrm C\)。
室内漏入热量速率
\(3.8\times10^4\,\mathrm{kJ/h}\)。
求制冷机至少需要多大的机械功率
（理想卡诺循环）。

\paragraph{习题 10.24}
可在严冬运行的“热泵”热机串联系统，
一台卡诺热机在热库 \(210^\circ\mathrm C\)
与河水 \(15^\circ\mathrm C\) 间工作，
其输出驱动一台逆卡诺制热机，
将河水中热量输送给取暖系统
\(60^\circ\mathrm C\) 的水。
设每燃烧 \(1\,\mathrm{kg}\) 煤释放
\(3.34\times10^7\,\mathrm J\)。
求取暖水的得热量及其相对燃煤热值的倍数。

\paragraph{习题 10.25}
图10.31致冷机循环
\(d\to a\to b\to c\to d\)；
\(d\to a\)、\(b\to c\) 绝热，
\(a\to b\) 在80℃等温压缩、液化，
\(c\to d\) 在5℃等温蒸发。
源教材状态表：
\[
\begin{array}{c|rrrrr}
 & T/^\circ\mathrm C&p/\mathrm{kPa}&V/\mathrm{m^3}&
 E/\mathrm{kJ}&\text{液态比例}/\%\\\hline
a&80&2305&0.0682&1969&0\\
b&80&2305&0.00946&1171&100\\
c&5&363&0.2202&1015&54\\
d&5&363&0.4513&1641&5
\end{array}
\]
（1）每循环从被冷却处吸热；
（2）向冷却器外空气排热；
（3）压缩机对工质做功；
（4）实际致冷系数，与两热库之间
理想卡诺致冷系数比较。

\paragraph{习题 10.26}
图10.32为逆向斯特林循环：
\(1\to2\) 高温 \(T_1\) 等温压缩，
\(3\to4\) 低温 \(T_2\) 等温膨胀，
\(2\to3\) 与 \(4\to1\) 为等容过程。
（1）证明用理想再生器时此循环致冷系数
与同温热库间逆卡诺循环相同；
（2）讨论若没有再生器，
在等容 \(4\to1\) 从高温热库吸热、
\(2\to3\) 把热量排回低温热库，
循环的有效致冷系数如何下降。
\textit{图示源题要求明确将等容热交换纳入
实际冷库净热量核算。}
